\section*{Capítulo \ref{ch:b1:rock-forming-minerals} --- Minerales formadores de rocas}

\begin{solution}{exo:b1:rock-forming-minerals:1}
\ce{SiO4^{4-}}: \omterm{def:b1:rock-forming-minerals:silicate-classes}{nesosilicato}, $s = 0$. \ce{Si2O7^{6-}}: \omterm{def:b1:rock-forming-minerals:silicate-classes}{sorosilicato}, $s = 1$. \ce{Si6O18^{12-}}:
\omterm{def:b1:rock-forming-minerals:silicate-classes}{ciclosilicato}, $s = 2$. \ce{Si2O5^{2-}}: \omterm{def:b1:rock-forming-minerals:silicate-classes}{filosilicato}, $s = 3$. \ce{SiO2}: \omterm{def:b1:rock-forming-minerals:silicate-classes}{tectosilicato}, $s = 4$.
\ce{SiO3^{2-}} tiene $4 - s/2 = 3$, así que $s = 2$: un anillo (\omterm{def:b1:rock-forming-minerals:silicate-classes}{ciclosilicato}) o una cadena simple
(\omterm{def:b1:rock-forming-minerals:silicate-classes}{inosilicato}); la razón por sí sola no permite distinguirlos.
\end{solution}

\begin{solution}{exo:b1:rock-forming-minerals:2}
Forsterita: $2\times24.305 + 28.085 + 4\times15.999 =
\qty{140.691}{g.mol^{-1}}$; fayalita: $2\times55.845 + 28.085 +
4\times15.999 = \qty{203.771}{g.mol^{-1}}$. Volúmenes molares:
$140.691/3.275 = \qty{42.96}{cm^3.mol^{-1}}$ y $203.771/4.392 =
\qty{46.40}{cm^3.mol^{-1}}$.
\end{solution}

\begin{solution}{exo:b1:rock-forming-minerals:3}
$(0.78 - 0.72)/0.72 = +\qty{8.3}{\percent}$ para \ce{Fe^{2+}};
$-\qty{4.2}{\percent}$ para \ce{Ni^{2+}}; $+\qty{15.3}{\percent}$ para
\ce{Mn^{2+}}; $+\qty{38.9}{\percent}$ para \ce{Ca^{2+}}. El hierro y el níquel, por debajo del
\qty{15}{\percent}, reemplazan libremente al magnesio; el manganeso, justo por encima, está limitado; el calcio,
más allá del \qty{30}{\percent}, apenas entra. Los cuatro tienen la misma carga que el magnesio.
\end{solution}

\begin{solution}{exo:b1:rock-forming-minerals:4}
$(3.51 - 2.16)/2.16 = \qty{62.5}{\percent}$ del grafito al diamante;
$(2.93 - 2.71)/2.71 = \qty{8.1}{\percent}$ de la calcita a la aragonita. Los \omterm{def:b1:rock-forming-minerals:polymorph}{polimorfos} más densos, el diamante y
la aragonita, son las formas de alta presión.
\end{solution}

\begin{solution}{exo:b1:rock-forming-minerals:5}
Con $x = 0.30$ en $\mathrm{Na}_{1-x}\mathrm{Ca}_x\mathrm{Al}_{1+x}\mathrm{Si}_{3-x}\mathrm{O}_8$:
$\mathrm{(Na_{0.7}Ca_{0.3})(Al_{1.3}Si_{2.7})O_8}$. Cargas:
$0.7 + 2\times0.3 + 3\times1.3 + 4\times2.7 = 0.7 + 0.6 + 3.9 + 10.8 = +16$,
y $8\times(-2) = -16$.
\end{solution}

\begin{solution}{exo:b1:rock-forming-minerals:6}
Una agrupación que se repite tiene dos tetraedros que comparten tres vértices (aquellos cuyo ápice une las dos
cadenas) y dos que comparten dos. Los cuatro tetraedros tienen $4\times4 = 16$ vértices, de los que $3 + 3 + 2 + 2 = 10$
son compartidos, es decir, $10/2 = 5$ oxígenos compartidos; los $16 - 10 = 6$ restantes no son compartidos. En total
$5 + 6 = 11$ oxígenos para 4 silicios: \ce{Si4O11}, con carga $4\times4 - 11\times2 = -6$. La proposición da lo mismo
con $s = 5/2$: $4 - 5/4 = 11/4$ oxígenos por silicio, carga $5/2 - 4 = -3/2$ por silicio.
\end{solution}

\begin{solution}{exo:b1:rock-forming-minerals:7}
Moles en \qty{100}{g}: \ce{SiO2} $40.20/60.083 = 0.66907$; \ce{MgO}
$47.10/40.304 = 1.16862$; \ce{FeO} $12.20/71.844 = 0.16981$; \ce{NiO}
$0.40/74.692 = 0.00536$. Oxígenos: $2\times0.66907 + 1.16862 + 0.16981 +
0.00536 = 2.68194$; $f = 4/2.68194 = 1.49146$. Por cuatro oxígenos: Si
\num{0.998}, Mg \num{1.743}, Fe \num{0.253}, Ni \num{0.008}. Las posiciones octaédricas contienen
$1.743 + 0.253 + 0.008 = 2.004$, la tetraédrica \num{0.998}: buen análisis,
\ce{(Mg_{1.743}Fe_{0.253}Ni_{0.008})SiO4}.
$\mathrm{Fo} = 100\times1.743/(1.743 + 0.253) = 87.3$. El factor $f$ multiplica Mg y Fe por igual, así que se cancela en
el cociente: $\mathrm{Fo} = 100\,(47.10/40.304)/(47.10/40.304 + 12.20/71.844) =
100\times1.16862/1.33843 = 87.3$.
\end{solution}

\begin{solution}{exo:b1:rock-forming-minerals:8}
Positivas: $1 + 2\times3 + 3\times4 + 3 = +22$; negativas:
$10\times(-2) + 2\times(-1) = -22$. Una lámina de composición \ce{Si4O10} llevaría $-4$
(\cref{prop:b1:rock-forming-minerals:si-o}, $s = 3$); sustituir un silicio de cada cuatro por aluminio añade una
carga negativa, así que \ce{(Si3Al)O10} lleva $-5$. Los dos aluminios octaédricos con los dos hidroxilos,
$6 - 2 = +4$, equilibran cuatro de las cinco cargas; el ion potasio entre las láminas equilibra la quinta.
\end{solution}

\begin{solution}{exo:b1:rock-forming-minerals:9}
(a) $M = (140.691 + 203.771)/2 = \qty{172.231}{g.mol^{-1}}$, $V = (42.96 +
46.40)/2 = \qty{44.68}{cm^3.mol^{-1}}$, $\rho = 172.231/44.68 =
\qty{3.855}{g.cm^{-3}}$; la media de las densidades es $(3.275 + 4.392)/2 =
\qty{3.834}{g.cm^{-3}}$, \qty{0.021}{g.cm^{-3}} menos. (b) Resuélvase
$203.771 + x(140.691 - 203.771) = 3.600\,\bigl(46.40 + x(42.96 - 46.40)\bigr)$:
$203.771 - 63.080x = 167.040 - 12.384x$, así que $x = 36.731/50.696 = 0.7245$: el
\omterm{def:b1:rock-forming-minerals:groups}{olivino} es $\text{Fo}_{72}$ (al entero por ciento más próximo).
\end{solution}

\begin{solution}{exo:b1:rock-forming-minerals:10}
Con $M_i = \rho_i V_i$,
\[
  \rho(x) - \bigl(x\rho_1 + (1 - x)\rho_2\bigr)
  = \frac{x\rho_1V_1 + (1 - x)\rho_2V_2 - \bigl(x\rho_1 + (1 - x)\rho_2\bigr)\bigl(xV_1 + (1 - x)V_2\bigr)}{V(x)} .
\]
Al desarrollar el producto, el numerador es $x\rho_1V_1 - x^2\rho_1V_1 +
(1 - x)\rho_2V_2 - (1 - x)^2\rho_2V_2 - x(1 - x)(\rho_1V_2 + \rho_2V_1)$, y
$x - x^2 = x(1 - x)$, $(1 - x) - (1 - x)^2 = x(1 - x)$, así que vale
$x(1 - x)(\rho_1V_1 + \rho_2V_2 - \rho_1V_2 - \rho_2V_1) =
x(1 - x)(\rho_1 - \rho_2)(V_1 - V_2)$. Para el \omterm{def:b1:rock-forming-minerals:groups}{olivino} (1 = forsterita) ambas
diferencias son negativas, $3.275 - 4.392 = -1.117$ y $42.96 - 46.40 =
-3.44$, así que su producto es positivo y la densidad queda por encima de la recta. En $\text{Fo}_{50}$:
$0.25\times1.117\times3.44/44.68 =
\qty{0.0215}{g.cm^{-3}}$, que concuerda con \cref{exo:b1:rock-forming-minerals:9} dentro del redondeo de los
volúmenes molares.
\end{solution}

\begin{solution}{exo:b1:rock-forming-minerals:11}
Total \qty{99.60}{\percent}. Moles de óxido: \ce{SiO2} $0.90042$,
\ce{Al2O3} $0.28344$, \ce{CaO} $0.19973$, \ce{Na2O} $0.08229$,
\ce{K2O} $0.00318$. Oxígenos: $1.80084 + 0.85033 + 0.19973 + 0.08229 +
0.00318 = 2.93636$; $f = 8/2.93636 = 2.72446$. Cationes (el doble de los moles de
\ce{Al2O3}, \ce{Na2O}, \ce{K2O}): Si \num{2.453}, Al \num{1.544}, Ca
\num{0.544}, Na \num{0.448}, K \num{0.017}. Posiciones tetraédricas:
$2.453 + 1.544 = 3.997$; posición grande: $0.544 + 0.448 + 0.017 = 1.009$. Entonces
$\mathrm{Al} - 1 = 0.544$, igual a Ca dentro del redondeo: cada calcio entró con un aluminio. $\mathrm{An} = 100\times0.544/1.009 = 53.9$:
una labradorita, $\text{An}_{54}$.
\end{solution}

\begin{solution}{exo:b1:rock-forming-minerals:12}
(a) $M(\ce{Fe3O4}) = 3\times55.845 + 4\times15.999 = \qty{231.531}{g.mol^{-1}}$.
Sus tres hierros, dados como \ce{FeO}, pesan $3\times71.844 = \qty{215.532}{g}$,
es decir, $215.532/231.531 = \qty{93.09}{\percent}$; el oxígeno ligado al
\ce{Fe^{3+}}, más allá de lo que da \ce{FeO}, falta en el total.
(b) $93.09/71.844 = 1.29572$ mol de \ce{FeO} y de oxígeno; $f =
4/1.29572 = 3.08708$; cuatro hierros por cuatro oxígenos, \ce{Fe4O4}. La magnetita tiene
tres cationes por cuatro oxígenos. Los oxígenos se contaron de menos: cada \ce{Fe^{3+}} tiene uno y medio
oxígenos, no uno como supone \ce{FeO}. (c) Sobre tres cationes, Fe $= 3$ con $a$ \ce{Fe^{2+}} y $b$ \ce{Fe^{3+}}:
$a + b = 3$ y $2a + 3b = 8$, así que $b = 2$, $a = 1$: $\mathrm{Fe^{2+}Fe^{3+}_2O_4}$. (d) Un \ce{FeO} y un
\ce{Fe2O3} por unidad de fórmula: $71.844/231.531 = \qty{31.03}{\percent}$ de \ce{FeO} y
$159.687/231.531 = \qty{68.97}{\percent}$ de \ce{Fe2O3}, total \qty{100.00}{\percent}.
\end{solution}

\begin{solution}{pb:b1:rock-forming-minerals:1}
\textbf{1.} \ce{SiO2} $28.085 + 2\times15.999 = 60.083$; \ce{FeO} 71.844;
\ce{MnO} 70.937; \ce{MgO} 40.304; \ce{NiO} 74.692; \ce{CaO} 56.077
(\si{g.mol^{-1}}).

\textbf{2.} $39.50 + 16.53 + 0.23 + 43.32 + 0.20 + 0.22 =
\qty{100.00}{\percent}$. Un total de \qty{97}{\percent} habría sugerido un componente no analizado (agua, u
oxígeno ligado a \ce{Fe^{3+}} dado como \ce{FeO}), o un análisis deficiente.

\textbf{3.} \ce{SiO2} $39.50/60.083 = 0.65742$; \ce{FeO} 0.23008; \ce{MnO}
0.00324; \ce{MgO} 1.07483; \ce{NiO} 0.00268; \ce{CaO} 0.00392.

\textbf{4.} $2\times0.65742 = 1.31485$ para \ce{SiO2}, los mismos números que en la
pregunta 3 para los demás; suma $1.31485 + 0.23008 + 0.00324 + 1.07483 +
0.00268 + 0.00392 = 2.62960$.

\textbf{5.} Cada óxido tiene un catión: los números de la pregunta 3.

\textbf{6.} $f = 4/2.62960 = 1.52114$.

\textbf{7.} Si $0.65742\times1.52114 = 1.000$; Fe 0.350; Mn 0.005; Mg 1.635;
Ni 0.004; Ca 0.006.

\textbf{8.} $1.635 + 0.350 + 0.005 + 0.004 + 0.006 = 2.000$ para las posiciones
octaédricas y 1.000 para la tetraédrica: exactamente el 2 y el 1 ideales.

\textbf{9.} \ce{(Mg_{1.635}Fe_{0.350}Mn_{0.005}Ni_{0.004}Ca_{0.006})SiO4}.

\textbf{10.} $2\times2.000 + 4\times1.000 = +8.000$ frente a $4\times(-2) =
-8$. Cada catión se contó con la carga de su óxido, así que las cargas positivas son el doble de los
oxígenos por construcción: la comprobación no puede fallar. Los totales por posición de la pregunta 8 sí
podrían haber fallado.

\textbf{11.} El níquel difiere del magnesio en
$(0.69 - 0.72)/0.72 = -\qty{4}{\percent}$, muy por debajo del \qty{15}{\percent}: sustituye libremente. El calcio
es $(1.00 - 0.72)/0.72 = \qty{39}{\percent}$ mayor, más allá del \qty{30}{\percent}: apenas entra.

\textbf{12.} El hierro contado como \ce{FeO} cuando parte de él es \ce{Fe^{3+}} da
demasiado pocos oxígenos por catión, de ahí más de tres cationes sobre cuatro oxígenos
(\cref{exo:b1:rock-forming-minerals:12}). La suma es $1.000 + 2.000 = 3.000$: sin \ce{Fe^{3+}} significativo.

\textbf{13.} $\mathrm{Fo} = 100\times1.635/(1.635 + 0.350) = 163.5/1.985 =
82.4$.

\textbf{14.} $\mathrm{Mg\#} = 82.4$: todo el hierro es \ce{Fe^{2+}} (pregunta 12), así que
$\mathrm{Fe^{2+}} = \mathrm{Fe}$ y las dos definiciones coinciden.

\textbf{15.} $\text{Fo}_{82.4}\text{Fa}_{17.6}$, $x = 0.824$; el hierro ocupa el
\qty{17.6}{\percent} de las posiciones octaédricas en el modelo binario
($0.350/2.000 = \qty{17.5}{\percent}$ con los elementos menores).

\textbf{16.} El factor multiplica Mg y Fe por igual y se cancela:
$\mathrm{Fo} = 100\times1.07483/(1.07483 + 0.23008) = 82.4$.

\textbf{17.} $40.5/40.304 = 1.00486$ y $20.0/71.844 = 0.27838$;
$\mathrm{Fo} = 100\times1.00486/1.28324 = 78.3$. El \omterm{def:b1:rock-forming-minerals:groups}{olivino} posterior es más rico en hierro.

\textbf{18.} $M = 140.691$ y \qty{203.771}{g.mol^{-1}}; sus volúmenes molares son
$V = 140.691/3.275 = 42.96$ y para el segundo, $203.771/4.392 =
\qty{46.40}{cm^3.mol^{-1}}$.

\textbf{19.} $M = 0.824\times140.691 + 0.176\times203.771 = 115.929 +
35.864 = \qty{151.793}{g.mol^{-1}}$; $V = 0.824\times42.96 +
0.176\times46.40 = 35.399 + 8.166 = \qty{43.565}{cm^3.mol^{-1}}$.

\textbf{20.} $\rho = 151.793/43.565 = \qty{3.484}{g.cm^{-3}}$.

\textbf{21.} $0.824\times3.275 + 0.176\times4.392 = 2.699 + 0.773 =
\qty{3.472}{g.cm^{-3}}$, \qty{0.012}{g.cm^{-3}} (\qty{0.3}{\percent}) por debajo del valor de mezcla ideal:
cerca, como predice \cref{exo:b1:rock-forming-minerals:10}.

\textbf{22.} $0.12\times\qty{3484}{kg.m^{-3}}\times\qty{1}{m^3} =
\qty{418}{kg}$.

\textbf{23.} $418/2950 = \qty{14.2}{\percent}$: el \omterm{def:b1:rock-forming-minerals:groups}{olivino} denso forma una parte mayor de la masa que del
volumen.

\textbf{24.} El \omterm{def:b1:rock-forming-minerals:groups}{olivino} es $\text{Fo}_{82}$: su contenido en forsterita es
\textbf{82 por ciento}.
\end{solution}
